\documentclass[11pt,a4paper]{article}
\usepackage[utf8]{inputenc}
\usepackage[T1]{fontenc}
\usepackage{lmodern}
\usepackage[margin=24mm]{geometry}
\usepackage{amsmath,amssymb,mathtools}
\usepackage{booktabs}
\usepackage{enumitem}
\usepackage{xcolor}
\usepackage[hidelinks]{hyperref}

\newcommand{\dd}{\mathop{}\!\mathrm{d}}
\newcommand{\ii}{\mathrm{i}}
\newcommand{\id}{\mathrm{id}}
\newcommand{\R}{\mathbb{R}}
\newcommand{\Z}{\mathbb{Z}}
\newcommand{\N}{\mathbb{N}}
\newcommand{\Ord}{\operatorname{ord}}
\newcommand{\Ker}{\operatorname{Ker}}
\newcommand{\Aut}{\operatorname{Aut}}
\newcommand{\slashed}[1]{\not\!#1}
\setlength{\parindent}{0pt}
\setlength{\parskip}{6pt}
\allowdisplaybreaks

\title{Particle Physics --- Strong Interaction}
\author{IMAPP lecture notes}
\date{Spring 2026}
\begin{document}\maketitle

\section{Deep inelastic scattering}
Deep inelastic scattering probes the internal structure of the proton at large momentum transfer. The relevant regime is characterized by
\[
Q^2=-q^2>0,
\]
with $Q^2$ large compared with hadronic scales. At sufficiently high energy, the proton behaves as a collection of point-like constituents, the partons.

\section{Quarks and colour}
Quarks come in six flavours and carry colour charge. The colour degree of freedom is described by the fundamental representation of $SU(3)$, with three colour states. Observable hadrons are colour singlets:
\[
3\otimes\bar 3=1\oplus8,\qquad 3\otimes3\otimes3=1\oplus\cdots.
\]
The absence of isolated coloured particles is confinement. The strong interaction is mediated by eight gluons, corresponding to the generators of $SU(3)$.

\section{Parton distribution functions}
The parton distribution function (PDF) $f_i(x,Q^2)$ gives the probability density to find a parton of type $i$ carrying momentum fraction $x$ at resolution $Q^2$. A schematic factorization formula is
\[
\sigma=\sum_i\int_0^1\dd x\,f_i(x,Q^2)\,\hat\sigma_i(x,Q^2),
\]
where $\hat\sigma_i$ is the hard partonic cross-section. PDFs are non-perturbative inputs and are measured or fitted from data.

\section{Hadron colliders}
At a collider such as the LHC, the partonic centre-of-mass energy is smaller than the proton-proton centre-of-mass energy because only fractions $x_1$ and $x_2$ of the proton momenta participate. The partonic energy satisfies $\hat s=x_1x_2s$. Increasing $Q^2$ resolves shorter distances and reveals more partonic structure.

\section{Comparison with the weak interaction}
The strong interaction is mediated by massless gluons and has a characteristic coupling that runs with the energy scale. The weak interaction is mediated by massive bosons, so its low-energy strength is suppressed by powers of $1/M_W^2$. Weak processes can change flavour and involve neutrinos; strong processes conserve flavour at the interaction vertex.

\section*{Summary}
Deep inelastic scattering established the partonic structure of the proton. QCD describes quarks carrying colour and interacting through eight gluons, while PDFs connect the non-perturbative proton to calculable hard scattering.
\end{document}
